\begin{table}[ht]
\centering
\caption{PTB-XL results, mean over $n$ random seeds. Each seed redraws the patient-level split and reinitialises the weights, so brackets give two-sided 95\% Student-$t$ intervals over that combined variation. The interval is a within-corpus statement and is not a generalisation interval.}
\label{tab:seed_ptbxl}
\small
\begin{tabular}{lcccc}
\toprule
Model & F1 & Prec. & Recall & $n$ \\
\midrule
InceptionTime & $0.706$ & $0.632$ & $0.825$ & 1 \\
ResNet1D & $0.702$ & $0.614$ & $0.844$ & 1 \\
SE-ResNet1D & $0.698$ & $0.620$ & $0.821$ & 1 \\
ConvNeXt-V2 & $0.691$ & $0.623$ & $0.785$ & 1 \\
CRNN & $0.682$ & $0.604$ & $0.806$ & 1 \\
Medformer & $0.674$ & $0.607$ & $0.758$ & 1 \\
Mamba & $0.672$ & $0.638$ & $0.711$ & 1 \\
\midrule
\multicolumn{5}{l}{\footnotesize Intervals describe how much each model moves when the split and the initialisation} \\
\multicolumn{5}{l}{\footnotesize are redrawn; they are not a pairwise test. See Table~\ref{tab:flip_survival}.} \\
\multicolumn{5}{l}{\footnotesize Averaging: macro-label.} \\
\bottomrule
\end{tabular}
\end{table}
